\section*{Wprowadzenie}
Źródła światła sztucznego przeszły w ciągu ostatniego wieku znaczną ewolucję. Od żarówek żarowych, przez lampy halogenowe i fluorescencyjne, po diody LED, każda technologia wnosi inne kompromisy między sprawnością, trwałością, kosztem i jakością barwy światła. Niniejszy przegląd omawia najważniejsze z nich.

\section*{Żarówki żarowe}
Światło powstaje w żarzącym się włóknie wolframowym, rozgrzanym do wysokiej temperatury przez przepływ prądu. Większość energii zamieniana jest na ciepło, a sprawność świetlna wynosi zwykle około 10–17 lm/W. Żywotność jest ograniczona, ponieważ włókno stopniowo paruje.

\section*{Żarówki halogenowe}
Wypełnienie bańki gazem halogenowym pozwala na regenerację parującego wolframu. Dzięki temu lampy mają dłuższą żywotność i bardziej stabilne parametry, przy zbliżonej sprawności do żarówek żarowych.

\section*{Świetlówki}
Źródłem światła jest wyładowanie elektryczne w parach rtęci, a luminofor na ściance bańki zamienia promieniowanie UV na widzialne. Sprawność jest wyraźnie wyższa niż w żarówkach żarowych, lecz lampy wymagają statecznika i zawierają rtęć, co komplikuje utylizację.

\section*{Diody LED}
Dioda półprzewodnikowa emituje światło bezpośrednio przy przepływie prądu. Współczesne modele osiągają sprawność rzędu kilkudziesięciu do ponad stu lm/W, a trwałość wynosi zwykle kilkadziesiąt tysięcy godzin. Parametrem wrażliwym pozostaje temperatura pracy, która wpływa na strumień światła i żywotność.

\section*{Porównanie}
\begin{tabularx}{\textwidth}{@{}lXXX@{}}
\toprule
\textbf{Technologia} & \textbf{Sprawność (lm/W)} & \textbf{Żywotność} & \textbf{Uwagi} \\
\midrule
Żarowa & ok. 10–17 & krótka & prosta konstrukcja \\
Halogenowa & ok. 15–25 & krótsza niż LED & gorąca bańka \\
Świetlówka & ok. 50–100 & średnia & zawiera rtęć \\
LED & ponad 100 & długa & wymaga odprowadzania ciepła \\
\bottomrule
\end{tabularx}

\section*{Wnioski}
Diody LED zdominowały rynek ze względu na sprawność i trwałość. Pozostałe technologie znajdują zastosowanie w niszowych obszarach, a wybór źródła światła zależy od wymagań dotyczących barwy, kosztu i warunków pracy.